% Copyright 2026 Allen Jue
%
% Licensed under the Apache License, Version 2.0 (the "License");
% you may not use this file except in compliance with the License.
% You may obtain a copy of the License at
%
%     http://www.apache.org/licenses/LICENSE-2.0
%
% Unless required by applicable law or agreed to in writing, software
% distributed under the License is distributed on an "AS IS" BASIS,
% WITHOUT WARRANTIES OR CONDITIONS OF ANY KIND, either express or implied.
% See the License for the specific language governing permissions and
% limitations under the License.
%
% kacho-article.tex — A realistic two-sided academic article with stock
% LaTeX margins, mathematics, tables, footnotes, and an inline \kachospray
% figure demonstrating zero text overlap.

\documentclass[11pt,a4paper,twoside]{article}
\usepackage{amsmath,amssymb}
\usepackage{booktabs}
\usepackage[
  seed=12,
  palette=hanga,
  theme=haru,
  layout=auto,
  density=0.65,
  birds=true,
  seal=true
]{kacho}

\title{\textbf{Algorithmic \emph{Kach\=o-e}: Procedural Botanical and Avian Illumination in Lua\LaTeX{} Margins}}
\author{Allen Jue\\ \small Studio for Computational Typography \& East Asian Brushwork}
\date{Spring 2026}

\begin{document}
\maketitle

\begin{abstract}
Traditional Japanese \emph{kach\=o-ga} (flower-and-bird paintings) and \emph{shin-hanga} woodblock prints achieve expressive naturalism through economy of gesture: a tapering ink-wash bough entering from a page edge, translucent mineral glazes pooled along petal rims, and a solitary songbird counterbalancing negative space (\emph{ma}). This article describes the mathematical and typographic architecture of \texttt{kacho}, a pure Lua\LaTeX{} package that synthesizes procedural watercolour and \emph{suibokuga} margin compositions directly at page shipout without external images or fonts.
\end{abstract}

\section{Introduction: Margin Illumination and Negative Space}

When illuminating a typeset page, ornament and prose share a fixed sheet of paper. Unlike static border clip-art, \emph{kach\=o-e} compositions are asymmetric and responsive to empty space. A plum branch (\emph{ume}) or cherry spray (\emph{sakura}) enters from the outer margin edge, arches toward the upper corner, and arrests its growth before touching a single glyph of the text block.\footnote{At shipout time, \texttt{kacho} traverses Lua\TeX's constructed \texttt{vlist} and \texttt{hlist} node trees to measure the exact bounding box of every glyph, rule, math list, and footnote.}

\begin{quote}
\itshape
Shiraume ya / sumi kanbashiki / K\=orokan\\[2pt]
White plum blossoms--- / and the fragrance of freshly ground ink / at the K\=orokan hall.\\[2pt]
\upshape\small --- Yosa Buson (1716--1784), \emph{Buson Kush\=u} (1784)
\end{quote}

Every page derives an independent deterministic pseudo-random stream from a document seed and the page index. Recompiling the same source always yields byte-identical vector paths, while consecutive pages draw from a shuffled composition deck so no two adjacent spreads repeat the same silhouette.

\section{Space Colonization and Pipe-Model Branching}

Woody sprays (\emph{ume}, \emph{sakura}, \emph{momiji}, \emph{matsu}, \emph{tsubaki}, and \emph{kinmokusei}) are grown inside the free margin envelope $\Omega \subset \mathbb{R}^2$ using a two-dimensional adaptation of the attractor-cloud space colonization algorithm of Runions, Lane, and Prusinkiewicz~(2007). Let $A = \{a_1, \dots, a_m\} \subset \Omega$ denote a Poisson-like sample of nutrient attractors inside the allowed margin band, and let $V = \{v_1, \dots, v_k\}$ denote the active skeleton vertices. At each growth step, each attractor $a_j$ influences its nearest skeleton node $v_i$ within radius $d_{\mathrm{reach}}$:
\begin{equation}
  \hat{n}_i = \frac{\sum_{a_j \in \mathcal{N}(v_i)} \frac{a_j - v_i}{\|a_j - v_i\|}}{\left\|\sum_{a_j \in \mathcal{N}(v_i)} \frac{a_j - v_i}{\|a_j - v_i\|}\right\|},
  \qquad
  v_{i}^{\mathrm{new}} = v_i + \delta \cdot \frac{\hat{n}_i + \tau \hat{g}}{\|\hat{n}_i + \tau \hat{g}\|},
  \label{eq:colonize}
\end{equation}
where $\delta$ is the internode step length and $\tau \hat{g}$ imparts seasonal phototropism or gravitropic droop (upward for gnarled winter plum, downward for cascading wisteria racemes).

Once the skeletal graph is complete, branch thicknesses are assigned from the twig tips back to the root following the botanical pipe model of Shinozaki et~al.~(1964). If a parent internode splits into child branches with radii $r_1, \dots, r_c$, cross-sectional vascular area is conserved according to
\begin{equation}
  r_{\mathrm{parent}}^p = \sum_{j=1}^{c} r_j^p, \qquad p \in [1.95, 2.15],
  \label{eq:pipemodel}
\end{equation}
producing thick gnarled basal boughs that taper organically into delicate terminal shoots.

\begin{table}[t]
\centering
\caption{Seasonal themes (\texttt{theme=}), primary botanical motifs, avian companions, and atmospheric accents in \texttt{kacho}.}
\label{tab:themes}
\medskip
\begin{tabular}{@{}lllll@{}}
\toprule
\textbf{Theme} & \textbf{Season} & \textbf{Primary Flora} & \textbf{Perched Bird} & \textbf{Atmosphere} \\
\midrule
\texttt{haru}        & Spring       & Sakura, Ume, Wisteria, Orchid     & Uguisu, Robin   & Kasumi mist \\
\texttt{natsu}       & Summer       & Iris (\emph{ayame}), Bamboo       & Uguisu, Swallow & Mt.\ Fuji, Bokashi \\
\texttt{aki}         & Autumn       & Momiji maple, Kiku, Osmanthus     & Mejiro, Sparrow & Ochre horizon \\
\texttt{fuyu}        & Winter       & Ume, Pine, Narcissus, Camellia    & Sparrow         & Scholar rock \\
\texttt{shochikubai} & New Year     & Pine, Bamboo, Plum (Trio)         & Uguisu          & Kinpaku gold \\
\bottomrule
\end{tabular}
\end{table}

\clearpage
\kachoThisPage{theme=natsu, layout=overhang, seed=23}

\section{Perceptual Pigment Mixing in Oklab}

Watercolour glazing requires smooth tonal transitions between concentrated mineral pigment and warm \emph{washi} paper. Linear interpolation in device RGB produces muddy grey midtones when mixing complementary pigments such as indigo (\emph{ai}) and cinnabar (\emph{shu}). Instead, \texttt{kacho} transforms linear sRGB tristimulus values $(R,G,B)$ into Bj\"orn Ottosson's perceptual Oklab space via cube-root cone responses:
\begin{equation}
  \begin{bmatrix} l \\ m \\ s \end{bmatrix} = \mathbf{M}_1 \begin{bmatrix} R \\ G \\ B \end{bmatrix},
  \qquad
  \begin{bmatrix} L \\ a \\ b \end{bmatrix} = \mathbf{M}_2 \begin{bmatrix} l^{1/3} \\ m^{1/3} \\ s^{1/3} \end{bmatrix}.
  \label{eq:oklab}
\end{equation}
All wet-on-wet washes (\emph{nijimi}), pooled back-run rims, and horizon gradations (\emph{bokashi}) interpolate in $(L,a,b)$ before converting back to sRGB PDF fill operators.

\begin{figure}[t]
\centering
\kachospray[width=0.72\linewidth, height=42mm, motif=ume, seed=19, from=left]
\caption{An inline botanical spray (\texttt{\textbackslash kachospray[width=0.72\textbackslash linewidth, height=42mm, motif=ume, seed=19, from=left]}) painted inside a reserved box in the normal figure flow. The layout engine automatically clips and bounds all strokes inside the box while keeping the caption untouched.}
\label{fig:inline-spray}
\end{figure}

As demonstrated in Figure~\ref{fig:inline-spray}, authors can also embed self-contained botanical studies directly inside floating figures or section dividers using \texttt{\textbackslash kachospray}. Each inline box records its elements with an \texttt{inline-} prefix and isolates its random stream from the surrounding page margin layout.\footnote{Changing the \texttt{seed} or \texttt{motif} of an inline \texttt{\textbackslash kachospray} never alters the outer margin composition of the page on which the float lands.}

\clearpage
\kachoThisPage{theme=aki, layout=corner-spray, seed=34}

\section{Calligraphic Brush Envelopes and Layered Washes}

A \emph{sumi-e} brushstroke is not a uniform-width polyline; the artist presses the belly of the goat-hair or wolf-hair brush onto the absorbent paper, pauses slightly at joints, and lifts the tip to a sharp point. Given an arc-length parameterized spine $\gamma(u) \in \mathbb{R}^2$ for $u \in [0,1]$ with unit normal $\hat{\nu}(u)$, \texttt{kacho} constructs closed left and right ribbon boundaries:
\begin{equation}
  \gamma_{\pm}(u) = \gamma(u) \pm \frac{1}{2} w_0 \, \phi(u) \Bigl(1 + \alpha \, \eta(f u + \theta_0)\Bigr) \hat{\nu}(u),
  \label{eq:ribbon}
\end{equation}
where $\phi(u)$ is the brush pressure profile (\texttt{taper}, \texttt{nail}, \texttt{bamboo}, or \texttt{leaf}), and $\eta(\cdot)$ is smoothstep-interpolated 1-D coherent noise that simulates hand tremor and paper tooth. When dry-brush \emph{kasure} is requested, the interior of the stroke is split into parallel bristle filaments separated by un-inked streaks of paper.

\begin{quote}
\itshape
Yama kurete / momiji no ake wo / ubaikeri\\[2pt]
As the mountain darkens, / twilight steals away the vermilion / of the maple leaves.\\[2pt]
\upshape\small --- Yosa Buson (1716--1784), \emph{Buson Kush\=u} (1784)
\end{quote}

Each petal and leaf is glazed in four passes inspired by Curtis et~al.~(1997) and adapted from the watercolour architecture of \texttt{florilegium}:
\begin{enumerate}
  \item \textbf{Wet body flood:} 2--4 jittered translucent polygon washes ($\alpha \approx 0.14$) establish the luminous base hue.
  \item \textbf{Multiply interior mottle:} A \texttt{/BM /Multiply} ExtGState pass deepens the petal base or calyx junction.
  \item \textbf{Back-run pooled rim:} A narrow darker perimeter stroke ($\alpha \approx 0.22$) reproduces capillary pigment migration toward drying wet edges.
  \item \textbf{Venation and lifted highlights:} Fine radiating veins (\texttt{vein\_fan}) and lifted dry-brush streaks reveal the underlying paper tone.
\end{enumerate}

\clearpage
\kachoThisPage{theme=fuyu, layout=rising, seed=45}

\section{Transactional Collision Avoidance at Shipout}

Safety of the typeset text is a hard invariant. Let $B_{\mathrm{paper}} = [0, W] \times [0, H]$ be the page rectangle, $B_{\mathrm{safe}}$ be the paper inset by the safety margin $s \ge 5\,\mathrm{mm}$, and $\{E_1, \dots, E_r\}$ be the union of the typearea and all measured glyph/rule/graphic bounding boxes inflated by the protective buffer $b \ge 4\,\mathrm{mm}$.

Every motif organ is drawn inside a transactional trial context (\texttt{ctx:try} or \texttt{Plan:attempt}). A candidate brushstroke or flower with bounding box $R$ is committed to the PDF stream if and only if
\begin{equation}
  R \subseteq B_{\mathrm{safe}}
  \qquad \text{and} \qquad
  \forall j \in \{1, \dots, r\}, \quad R \cap E_j = \emptyset.
  \label{eq:safety}
\end{equation}
If a candidate branch or blossom violates~\eqref{eq:safety}, the trial buffer is discarded without emitting a single PDF operator, and the layout engine either shrinks the scale factor $s \in \{1.0, 0.80, 0.64, 0.50\}$ or prunes the offending twig. Consequently, blossoms never get sliced in half by a clipping rectangle, and body text remains completely untouched.

\begin{quote}
\itshape
Iza saraba / yukimi ni korobu / tokoro made\\[2pt]
Come then, let us go / out to view the falling snow--- / till we tumble over!\\[2pt]
\upshape\small --- Matsuo Bash\=o (1644--1694), \emph{Oi no Kobumi} (1688)
\end{quote}

\section{Conclusion}

By uniting Lua\TeX{} shipout box inspection, space-colonization branching, pipe-model allometry, and Oklab watercolour glazing, \texttt{kacho} brings the quiet elegance of Japanese \emph{kach\=o-e} album leaves to standard \LaTeX{} documents with zero external asset dependencies.

\end{document}
